\documentclass[a0paper,portrait]{baposter}
\usepackage[french]{babel}


\usepackage{wrapfig}
\usepackage{lmodern}
\usepackage{lipsum,graphicx}
\usepackage{fontspec}
\usepackage{advdate}
\usepackage{svg}
\usepackage{xcolor}
\usepackage{tabularx}


% Set main font (Abraham)
%\setmainfont[
%    Path = font/Abraham/,
%    Extension = .otf,
%    % SmallCapsFont = *-smallcaps,
%    BoldFont = *-bold, % Outline
%    ItalicFont = *-italic, % Italic
%    BoldItalicFont = *-bolditalic % Outline Italic
%]{abraham}

%

\graphicspath{{figures/}} % Directory in which figures are stored

\newcommand{\compresslist}{%
\setlength{\itemsep}{0pt}%
\setlength{\parskip}{1pt}%
\setlength{\parsep}{0pt}%
}

\newenvironment{boenumerate}
  {\begin{enumerate}\renewcommand\labelenumi{\textbf\theenumi.}}
  {\end{enumerate}}

\usepackage{enumitem}
\setlist[1]{leftmargin=2em}
\setlist[2]{leftmargin=2em}

\begin{document}

\definecolor{Mycolor1}{HTML}{3FFF00}
\definecolor{Mycolor2}{HTML}{0FF000} %0FF000

\begin{poster}
{
grid=false,
headerborder=open, % Adds a border around the header of content boxes
colspacing=1em, % Column spacing
bgColorOne=white, % Background color for the gradient on the left side of the poster
bgColorTwo=white, % Background color for the gradient on the right side of the poster
borderColor=Mycolor1, % Border color
headerColorOne=Mycolor2, % Background color for the header in the content boxes (left side)
headerColorTwo=Mycolor2, % Background color for the header in the content boxes (right side)
headerFontColor=white, % Text color for the header text in the content boxes
boxColorOne=white, % Background color of the content boxes
textborder=rounded, %rectangle, % Format of the border around content boxes, can be: none, bars, coils, triangles, rectangle, rounded, roundedsmall, roundedright or faded
eyecatcher=false, % Set to false for ignoring the left logo in the title and move the title left
headerheight=0.08\textheight, % Height of the header
headershape=rounded, % Specify the rounded corner in the content box headers, can be: rectangle, small-rounded, roundedright, roundedleft or rounded
headershade=plain,
headerfont=\Large\textsf, % Large, bold and sans serif font in the headers of content boxes
%textfont={\setlength{\parindent}{1.5em}}, % Uncomment for paragraph indentation
linewidth=1pt % Width of the border lines around content boxes
}
{}
%
%----------------------------------------------------------------------------------------
%	TITLE AND AUTHOR NAME
%----------------------------------------------------------------------------------------
%
{\textsf{{Tableau de bord - Éco-marathon Shell 2023}}} % Titre en Sans Serif
{\sf\vspace{0.1em}\\
Électrique  -  {\AdvanceDate[0]\today} 
\vspace{0.1em}\\
\small{ Période couverte du {\AdvanceDate[-9]\today} au {\AdvanceDate[-3]\today}
}
}
{\includegraphics[width=.1\linewidth]{img/udes.pdf}} % TU Dublin logo


% this states the box starts at column 0 (edge of page), row 0 (top of page) for a span of 3 (columns wide)
\headerbox{Objectif de projet}{name=objectif,column=0,row=0, span=3}{
Créer un véhicule efficace à propulsion électrique pour participer et gagner la compétition Shell Eco-marathon dans la division ``véhicule urbain'' en 2023.\\[15pt]
\textbf{QUOTE} : On veut le beurre et l'argent du beurre.
}


% this states the box starts at column 0 (edge of page), directly below the box labelled introduction for a span of 1 (column wide)
\headerbox{Objectifs de session}{name=objectif_session,column=1,below=objectif,span=2}{

%\vspace{0.15cm}
\centering
{    \begin{tabularx}{\linewidth}{
        >{\hsize=1.5\hsize}X
        >{\hsize=1.0\hsize}X
        >{\centering\arraybackslash\hsize=0.5\hsize}X
      }
        
        \textbf{Objectif} & \textbf{Système} & \textbf{Avancement} \\
         Modélisation numérique & Moteur/Transmission & 100\% \\
         CAD détaillé & Moteur/Transmission & 99\% \\
         Assemblage V0 & Moteur/Transmission & 32\% \\
         PCB V2 onduleur & Onduleur & 75\% \\
         PCB de contrôle fonctionnel & Instrumentation & 30\% \\
         BMS V2 commandé & Bloc batterie & 60\% \\
         Bloc batterie V0 assemblé et testé & Bloc batterie & 85\% \\
         PCBs de lumière fonctionnels et assemblés & Instrumentation & 60\% \\
         Essuie & Instrumentation & 40\% \\
      \end{tabularx}
    
        
     % Template des Objectifs rencontés:
     %
     % Objectif : l'objectif de la session
     % Système  : le nom ou numéro système que l'objectif touche ex: Simulateur ou SIM1
     % &  & = ESPACE POUR LE FORMATAGE NE PAS ENLEVER
     % Avancement : Pourcentage estimé d'avancement de l'objectif.
     %
     %  
     %  \textbf{Objectif} & \textbf{Système} &  & \textbf{Avancement}\\
     %   Objectif & Système &  & Avancement \\
     %   Objectif & Système &  & Avancement \\
     %   Objectif & Système &  & Avancement \\
     %   Objectif & Système &  & Avancement \\
     %   Objectif & Système &  & Avancement \\
    
    
    
    
    
}

%\vspace{0.2cm} % spacing vertical

}

% this states the box starts at column 0 (edge of page), directly below the box labelled subtopic1 for a span of 1 (column wide)
\headerbox{Ordre du jour}{name=ordre_jour,column=0,below=objectif,span=1}{

\begin{itemize}
    \item \textbf{Rencontre gestion :}
    \begin{enumerate}
        \item Red flags 
        \item Tour de table sur les avancements des objectifs
    \end{enumerate}
    \item \textbf{Rencontre technique :}
    \begin{enumerate}
        \item Overview complet des systèmes
    \end{enumerate}
\end{itemize}

\vspace{0.05cm} % spacing vertical
}

\headerbox{Membres}{name=membres,column=2,below=ordre_jour,span=1}{
\vspace{0.3cm}


% inserts an image inside the box, 5 rows high | r = on the right along side text 
\begin{wrapfigure}[5]{r}{0.25\textwidth}
\includegraphics[width=0.9\linewidth]{img/membres/François-Paquette-3.jpg} 
\end{wrapfigure}
\subsubsection*{}
\vspace{-2mm}
\textbf{François Paquette}

Batterie / BMS

Moteur

\begin{wrapfigure}[5]{r}{0.25\textwidth}
\includegraphics[width=.9\linewidth]{img/membres/Gabriel-Cabana-2.jpg} 
\end{wrapfigure}
\subsubsection*{}
\vspace{-2mm}
\textbf{Gabriel Cabana}

Moteur

Instrumentation


\begin{wrapfigure}[5]{r}{0.25\textwidth}
\includegraphics[width=.9\linewidth]{img/membres/Jérôme-Gelé-2.jpg} 
\end{wrapfigure}
\subsubsection*{}
\vspace{2mm}
\textbf{Jérôme Gelé}

Contrôle moteur \& Onduleur

\begin{wrapfigure}[5]{r}{0.25\textwidth}
\includegraphics[width=.9\linewidth]{img/membres/Joël-Grégoire-Lagueux-1.jpg} 
\end{wrapfigure}
\subsubsection*{}
\vspace{2mm}
\textbf{Joël Grégoire-Lagueux}

Instrumentation

\begin{wrapfigure}[5]{r}{0.25\textwidth}
\includegraphics[width=.9\linewidth]{img/membres/Loïc-Poirier-2.jpg} 
\end{wrapfigure}
\subsubsection*{}
\vspace{-2mm}
\textbf{Loïc Poirier}

Batterie / BMS

Simulateur

\begin{wrapfigure}[5]{r}{0.25\textwidth}
\includegraphics[width=.9\linewidth]{img/membres/Marc-Antoine-Dubreuil-2.jpg} 
\end{wrapfigure}
\subsubsection*{}
\vspace{-2mm}
\textbf{Marc-Antoine Dubreuil}

Contrôle moteur \& Onduleur

Simulateur

\begin{wrapfigure}[5]{r}{0.25\textwidth}
\includegraphics[width=.9\linewidth]{img/membres/Thomas-Chagnon-3.jpg} 
\end{wrapfigure}
\subsubsection*{}
\vspace{2mm}
\textbf{Thomas Chagnon}

Moteur / Transmission

\begin{wrapfigure}[5]{r}{0.25\textwidth}
\includegraphics[width=.9\linewidth]{img/membres/Vincent-Bonneau-3.jpg} 
\end{wrapfigure}
\subsubsection*{}
\vspace{2mm}
\textbf{Vincent Bonneau}

Contrôle moteur \& Onduleur


\begin{wrapfigure}[5]{r}{0.25\textwidth}
\includegraphics[width=.9\linewidth]{img/membres/Xavier-Morin-2.jpg} 
\end{wrapfigure}
\subsubsection*{}
\vspace{2mm}
\textbf{Xavier Morin}

Moteur / Transmission

Transferts thermiques


\vspace{1cm}
%\begin{wrapfigure}[8]{r}{0.25\textwidth}
%\includegraphics[width=.9\linewidth]{DevOps.png} 
%\includegraphics[width=.9\linewidth]{DevOps.png} 
%\end{wrapfigure}
%
%\vspace{0.3cm}
%\subsubsection*{2. Another sub section:}
%
%Subsection description
%
%\vspace{2.3cm} %remove this, only added for spacing
%
%\begin{wrapfigure}[5]{r}{0.25\textwidth}
%\includegraphics[width=.9\linewidth]{TUDublin.jpg} 
%\end{wrapfigure}
%
%\vspace{0.35cm}
%\subsubsection*{3. Another sub section:}
%Final Subsection description
%
%\vspace{2cm} %remove this, only added for spacing
}


\headerbox{Courbe en S}{name=courbe_s,column=0,span=2,below=ordre_jour}{
{
    {
        \hspace{0.75cm}
        \centering
    	\includegraphics[width=0.9\linewidth]{img/Courbe_S.pdf}
    }
}

}

\headerbox{Heures travaillées}{name=heures,column=0,span=2,below=courbe_s}{
{
    
    {
        \hspace{0.75cm}
        \centering
    	\includegraphics[width=0.9\linewidth]{img/Heures_travaillees.pdf}
    }
    % \vspace{1.3cm}
}

}

\end{poster}

\begin{poster}
{
grid=false,
headerborder=open, % Adds a border around the header of content boxes
colspacing=1em, % Column spacing
bgColorOne=white, % Background color for the gradient on the left side of the poster
bgColorTwo=white, % Background color for the gradient on the right side of the poster
borderColor=Mycolor1, % Border color
headerColorOne=Mycolor2, % Background color for the header in the content boxes (left side)
headerColorTwo=Mycolor2, % Background color for the header in the content boxes (right side)
headerFontColor=white, % Text color for the header text in the content boxes
boxColorOne=white, % Background color of the content boxes
textborder=rounded, %rectangle, % Format of the border around content boxes, can be: none, bars, coils, triangles, rectangle, rounded, roundedsmall, roundedright or faded
eyecatcher=false, % Set to false for ignoring the left logo in the title and move the title left
headerheight=0\textheight, % Height of the header
headershape=rounded, % Specify the rounded corner in the content box headers, can be: rectangle, small-rounded, roundedright, roundedleft or rounded
headershade=plain,
headerfont=\Large\textsf, % Large, bold and sans serif font in the headers of content boxes
%textfont={\setlength{\parindent}{1.5em}}, % Uncomment for paragraph indentation
linewidth=1pt % Width of the border lines around content boxes
}
{}
{\textsf{{}}} % Titre vide pcq sinon ça plante


\headerbox{Tâches accomplies}{name=taches,column=0,row=0, span=3}{
{\begin{tabularx}{\linewidth}{
    |>{\hsize=2.25\hsize}X|% 10% of 4\hsize 
    >{\hsize=0.75\hsize}X|% 30% of 4\hsize
    >{\hsize=0.75\hsize}X|% 30% of 4\hsize 
    >{\centering\arraybackslash\hsize=0.25\hsize}X|% 30% of 4\hsize 
       % sum=0.2\hsize for 4 columns
  }
    \hline
    \textbf{Tâche} & \textbf{Système} & \textbf{Responsable} & \textbf{Heures}\\\hline
     Modélisation & Thermique & Xavier Morin & 4\\\hline
     Modélisation et calculs & Transmission & Xavier Morin & 9\\\hline
     Steering Committee & Gestion & Xavier Morin & 2\\\hline
     Choix du circuit intégré du système de gestion des batteries & Bloc batterie / BMS & Loïc Poirier & 4\\\hline
     Début de la conception du convertisseur CC/CC & Bloc batterie / BMS & Vincent Bonneau & 4\\\hline
     Lire la documentation sur le contrôle moteur & Contrôle moteur & Jérôme Gelé & 2\\\hline
     Choix du circuit intégré du système de gestion des batteries & BMS & Jérôme Gelé & 2\\\hline
     Conception préliminaire du circuit & Bloc batterie / BMS & Loïc Poirier & 4\\\hline
     Modélisation du contrôle moteur & Contrôle moteur & Vincent Bonneau & 4\\\hline
    
     Schématique de l'onduleur et librairies Altium & Contrôle moteur & Marc-Antoine Dubreuil & 6\\\hline
     Modélisation préliminaire du moteur en accélération & Moteur & Gabriel Cabana & 12\\\hline
  \end{tabularx}

% TEMPLATE des lignes du tableau de taches accomplis cette semaine
%
% Tâche: résumé de la tâche
% Système : Numéro de système ou nom complet ex : Simulateur ou SIM1
% Responsable : Initiales ou nom complet ex : Gabriel Cabana ou G.C.
% Heures : Heures passées par le responsable à faire la tâche la semaine dernière (jeudi à jeudi)
%
%\textbf{Tâche} & \textbf{Système} & \textbf{Responsable} & \textbf{Heures}\\\hline
%    Tâche & Système & Responsable & Heures\\\hline
%    Tâche & Système & Responsable & Heures\\\hline
%    Tâche & Système & Responsable & Heures\\\hline
%    Tâche & Système & Responsable & Heures\\\hline
%    Tâche & Système & Responsable & Heures\\\hline
%    Tâche & Système & Responsable & Heures\\\hline
%    Tâche & Système & Responsable & Heures\\\hline
%    Tâche & Système & Responsable & Heures\\\hline
%    Tâche & Système & Responsable & Heures\\\hline
%    Tâche & Système & Responsable & Heures\\\hline
%    Tâche & Système & Responsable & Heures\\\hline
%    Tâche & Système & Responsable & Heures\\\hline
%  }

%\vspace{2cm} %remove this, only added for spacing

}

% this states the box starts at column 0 (edge of page), directly below the box labelled subtopic1 for a span of 1 (column wide)

\headerbox{Problèmes}{name=problemes,column=0,below=taches,span=3}{

{\begin{tabularx}{\linewidth}{
    |>{\hsize=1.7\hsize}X|% 10% of 3\hsize 
    >{\hsize=0.5\hsize}X|% 30% of 3\hsize
    >{\hsize=0.8\hsize}X|% 30% of 3\hsize 
       % sum=0.2\hsize for 4 columns
  }
    \hline
    \textbf{Problème} & \textbf{Système} & \textbf{Responsable} \\\hline
    Quincaillerie manquante J.M. Bergeron & Transmission & Xavier Morin et Thomas Chagnon\\\hline
    Assemblage soudé arbres des roues & Transmission & Xavier Morin et Thomas Chagnon\\\hline
    Assemblage logement à roulements et longeron véhicule & Transmission & Xavier Morin et Thomas Chagnon\\\hline
    Mauvaise organisation de la documentation PCB de contrôle & Instrumentation & Vincent Bonneau\\\hline
\end{tabularx}
    
    
 % Template des problèmes rencontés:
 %
 % Problème : le problème rencontré cette semaine
 % Système  : le nom ou numéro système que le problème touche ex: Simulateur ou SIM1
 % Responsable : Nom ou Initiales du ou des personnes touchées par ce problèmes ex: G.C. C.E.G.
 %
 %  \hline
 %  \textbf{Problème} & \textbf{Système} & \textbf{Responsable}\\\hline
 %  Problème & Système & Responsable \\\hline
 %  Problème & Système & Responsable \\\hline
 %  Problème & Système & Responsable \\\hline
 %  Problème & Système & Responsable \\\hline
 %  Problème & Système & Responsable \\\hline
 %  Problème & Système & Responsable \\\hline
 %  Problème & Système & Responsable \\\hline
 %}

%\vspace{.65cm}
%\begin{itemize}
%    \item \textbf{Faire un schéma de dataflow du simulateur:} Mauvais compréhension du flow de données du %simulateur dans l'équipe informatique				
%    \item \textbf{Peux pas entrer les heures dans le DVP:} Courbe en S est pas représentative.	
%    \end{itemize}%

%    \vspace{1.8cm} % spacing vertical

}

% this states the box starts at column 0 (edge of page), directly below the box labelled introduction for a span of 1 (column wide)
\headerbox{Avancement systèmes}{name=systemes,column=0,below=problemes,span=3}{

%\vspace{0.15cm}
\centering
{\includegraphics[width=1\linewidth]{img/avancement.pdf}}

%\vspace{0.2cm} % spacing vertical

}



\headerbox{Risques}{name=risques,column=0,span=3,below=systemes}{
{
    {\begin{tabularx}{\linewidth}{
    |>{\let\newline\\\hsize=0.40\hsize}X|% 10% of 4\hsize 
    >{\hsize=0.25\hsize}X|% 30% of 4\hsize
    >{\hsize=0.25\hsize}X|% 30% of 4\hsize 
    >{\centering\arraybackslash\hsize=0.1\hsize}X|% 30% of 4\hsize 
       % sum=0.2\hsize for 4 columns
  }
    \hline
    \textbf{Risque} & \textbf{Prévention} & \textbf{Gestion} & \textbf{Priorité}\\\hline
  % Mauvaises connexion des phases du moteur & revues des pairs avant exécution & Rebobinage du moteur ou utilisation du moteur de remplacement & HAUTE\\\hline
  
    Moteur V0 non fonctionnel & Mettre les bouchées double afin d'avoir un prototype & Se limiter à un moteur acheté ou faire usiner à l'extérieur & HAUTE\\\hline
    
    Moteur V0 ou acheté non caractérisé pour l'échéancier du 6 octobre & Intégration partielle au véhicule & Communication active avec les équipes interdépendantes & MOYENNE\\\hline
    
    %Moteur qui ne respecte pas les spécifications acceptables & Faire simulations, calculs, revues de conception avec experts & Achat d'un moteur plan B avec date limite & HAUTE\\\hline
    

    Mauvaise conception système onduleur & Revue de conception avec experts, validation V1, validation PCB contrôle, attribuer MVP & acquisition onduleur (att. Budget/Temps) & HAUTE\\\hline
    
    Mauvaise conception du PCB de lumière & Système simple et revue de conception inter-équipe & Modifications directes sur le PCB ou achat de composants aditionnels au besoin & HAUTE\\\hline
    
    Commande des connecteurs tardive & Communication directe avec commanditaire & Commande faite par l'équipe au lieu du commanditaire & MOYENNE\\\hline
    

  \end{tabularx}
  
  
% Template pour le tableau des risques de la semaine : 

% Risque : le risque de la semaine
% Mitigation: comment mitiger le risque pour réduire ses conséquences/occurances/etc.
% Conséquence : La ou les conséquences si le risque survient
% Priorité : Priorité du risque sur une échelle de 1 à 5.  5 étant + prioritaire. La priorité est basé sur les conséquences la probablité d'occurence, etc.  
  
%  \textbf{Risque} & \textbf{Mitigation} & \textbf{Conséquence} & \textbf{Priorité}\\\hline
%    Risque & Mitigation & Conséquence & Priorité\\\hline
%    Risque & Mitigation & Conséquence & Priorité\\\hline
%    Risque & Mitigation & Conséquence & Priorité\\\hline
%    Risque & Mitigation & Conséquence & Priorité\\\hline
%    Risque & Mitigation & Conséquence & Priorité\\\hline
%    Risque & Mitigation & Conséquence & Priorité\\\hline
%    Risque & Mitigation & Conséquence & Priorité\\\hline
%    Risque & Mitigation & Conséquence & Priorité\\\hline
%    Risque & Mitigation & Conséquence & Priorité\\\hline
%    Risque & Mitigation & Conséquence & Priorité\\\hline
%    Risque & Mitigation & Conséquence & Priorité\\\hline
%    Risque & Mitigation & Conséquence & Priorité\\\hline}
}

}

\headerbox{Objectifs semaine prochaine}{name=objectifs_prochain,column=2,span=1,below=risques}{
{
    % \vspace{.65cm}
    \begin{itemize}
    \item \textbf{Moteur et Transmission :}
    \begin{enumerate}
        \item Terminer l'analyse analytique des efforts dans les shafts de transmission pour pouvoir dimensionner les shafts correctement.
        \item Modifier les pièces en fonction du retour de la revue de conception.
    \end{enumerate}
    \item \textbf{Contrôle moteur et Onduleur :}
    \begin{enumerate}
        \item Finalisation du design de l'onduleur
        \item Support à l'équipe informatique
    \end{enumerate}
    \item \textbf{Batterie BMS et Instrumentation :}
    \begin{enumerate}
        \item Finaliser le design des convertisseurs
        \item Révison des pièces et du circuit du BMS
    \end{enumerate}
    \item \textbf{Audit d'équipe :}
    \begin{enumerate}
        \item Compléter budget par système
        \item Faire présentation d'équipe
    \end{enumerate}
    \end{itemize}
    
    % \vspace{1cm} % spacing vertical
}

}

\headerbox{Cashflow}{name=cashflow,column=0,span=2,below=risques}{
{

    % \vspace{.6cm} % spacing vertical
    {
        \hspace{1cm}
        \centering
    	\includegraphics[width=0.85\linewidth]{img/cashflow.pdf}
    }
    %\vspace{1.15cm} % spacing vertical
}

}

\end{poster}




\end{document}
